\section{Related Work}\label{sec:related}

\subsection{Text-to-CAD: from command sequences to executable programs}
\label{sec:related:t2c}

Learning-based CAD generation was first formulated as sequence modeling:
DeepCAD represents models as sketch-and-extrude command sequences and trains
a generative network over them \cite{deepcad}, and Text2CAD pairs such
sequences with natural-language descriptions at multiple levels of expertise
\cite{text2cad}. The shift to large language models recast CAD generation as
\emph{program synthesis}: LLM4CAD emits a Python-based domain-specific
language \cite{llm4cad}, Text-to-CAD aligns parametric data with language via
a vision--language model \cite{texttocad_vl}, and CAD-GPT augments multimodal
LLMs with spatial-reasoning prompts to synthesize construction sequences
\cite{cadgpt}. More recent systems generate directly executable CadQuery
programs and run them to obtain B-rep geometry \cite{procad,cadrille}.

Executable Python programs rather than opaque token vocabularies are now the
preferred target representation as well: CADRecode reverse-engineers point
clouds directly into CadQuery scripts with a small LLM and a
million-program synthetic corpus \cite{cadrecode}, and CADFS collects a
million-scale corpus of \emph{real} FeatureScript programs with complete
modeling histories, executes them through a CAD kernel to obtain B-reps, and
evaluates on Chamfer, normal consistency, and edge Chamfer distance
\cite{cadfs}; CAD-Llama adapts open-weight LLMs to the same
program-synthesis setting \cite{cadllama}. We adopt this executable-program
view---and CADFS's multi-metric evaluation, including edge-aware distances
that expose fine structural defects---but generate forward from language
while tracking where every required value came from.

Two further trends are directly relevant. First, \emph{structured
intermediate tokens help}: FlexCAD achieves controllable generation and
partial editing through hierarchy-aware masking of structured CAD text
\cite{flexcad}, and Pointer-CAD unifies B-rep and command sequences with
explicit pointer-based edge and face references \cite{pointercad}---structure
in the representation buys controllability that raw free text cannot. Second,
\emph{visual and geometric feedback during training helps the generator}:
CADFusion alternates between learning a visual-feedback model and fine-tuning
the LLM on that feedback \cite{cadfusion}, while FutureCAD grounds generated
programs onto text-described B-rep primitives and uses the grounding signal as
a reinforcement reward \cite{futurecad}. Crucially, both use feedback as a
\emph{training-time} signal over the generator; at inference time the user's
specification is still consumed once, as free text, with no record of which
values were stated, assumed, or corrected.

Multimodal unified representations such as CAD-MLLM \cite{cadmllm} and
general-purpose CAD agents that combine tools with vision--language models
\cite{caddesigner,cadassistant} broaden the range of solvable tasks, but do
not address the provenance and verification of the design requirements
themselves---the gap we target.

\subsection{Proactive clarification of user intent}\label{sec:related:clarify}

Rather than forcing a single prompt to carry every requirement, proactive
agents ask before acting; the idea appears in task-oriented dialogue and
information-seeking systems \cite{clarify_general} and, in CAD, is
instantiated by ProCAD \cite{procad}: a Clarifier detects underspecified or
conflicting prompts and asks targeted questions before a separately trained
Coder produces CadQuery. Misleading prompts---parameters removed
(underspecification) or duplicated with contradictory values (direct
conflict)---are synthesized at controlled severity $K$, and an oracle user
simulator answers from the ground-truth description. We build directly on
this problem formulation and reuse its benchmark protocol.

Our analysis, however, identifies a structural weakness: the two agents
exchange a single rewritten description with no schema, no per-parameter
provenance, and no post-generation check. Information lost in that rewrite is
unrecoverable downstream, and a wrong value carried through clarification is
undetectable without the ground-truth mesh that the benchmark happens to
provide. We replace the free-text hand-off with a provenance-tracked
specification, make the answering user explicitly imperfect in our
experiments, and add a verification loop the original pipeline lacks.

\subsection{Structured and editable intermediate representations}
\label{sec:related:structured}

Exposing intermediate structure has a long history in procedural and
constructive-solid-modeling systems, where CSG trees and feature histories
make geometry inspectable and locally editable
\cite{requicha,featurescript,openscad}. The recent generative literature
returns to this idea: CAD-Editor locates a target span in CAD text and
re-infills only that span, avoiding full re-generation for edits
\cite{cadeditor}; FlexCAD's hierarchy masking similarly preserves non-target
parts \cite{flexcad}.

Most directly related is Treellis \cite{treellis}, which has an LLM emit an
assembly modeling language (AML)---a CSG-style tree of primitives with
editable parameters---and presents the tree, viewport, and parameter panel as
linked, editable views; a scoped edit sends only the selected subtree to the
model. In a controlled study, structured editing largely eliminated the
``regenerate everything and watch unrelated parts change'' loop that plain
prompting produces. We take this as strong motivation that users and agents
alike need an intermediate object between language and geometry.

But Treellis's own measurements expose the other half of the problem: only 32
of 50 generated trees were editable on the first pass (43 with exemplars and
structural constraints), because nothing verifies that the tree faithfully
captures the intent or that its parameters are complete. Their tree is a
\emph{geometric} structure revealed after generation for a human to repair;
our specification is an \emph{intent} structure built before generation, with
every value attributed to a source, and checked by the agent itself without a
human in the loop. The two are complementary downstream: a
provenance-tagged specification is a natural authoring layer for an editable
assembly tree.

\subsection{Verifying and repairing generated programs}
\label{sec:related:verify}

Self-verification of LLM outputs ranges from consistency sampling
\cite{selfcheckgpt} and self-refinement \cite{selfrefine} to execution-based
repair, where compiler and runtime feedback drives another generation
attempt; in CAD, cadrille trains with online RL against execution feedback
\cite{cadrille}. The closest system, CADCodeVerify \cite{cadcodeverify},
prompts a VLM to (i) generate binary validation questions from the
description, (ii) answer them by looking at multi-view renders of its own
generated object, and (iii) rewrite code from the resulting feedback, with a
separate repair loop for compilation failures.

This verifies geometry \emph{visually and self-referentially}: the questions,
the answers, and the judged render all come from the same model, with no
grounded reference, and holistic image questions are insensitive to local
defects. SEEK-CAD applies the same render-and-critique loop with
chain-of-thought reasoning on open-weight models \cite{seekcad}, and ReCAD
improves the generator with reinforcement learning using visual feedback as
reward \cite{recad}; both inherit the unanchored, holistic nature of the
visual signal.

Closest to our goals is CADDesigner \cite{caddesigner}, a general-purpose
ReAct-style agent for conceptual CAD that combines three ingredients we also
rely on: a \emph{requirement-expansion} subagent that turns a vague request
into dimensions, constraints, assumptions, APIs, and an ordered modeling
process; an explicit-context imperative paradigm (ECIP) that keeps modeling
context and operation intent explicit during code generation; and iterative
repair driven by execution tracebacks and rendered visual feedback. It is,
to our knowledge, the clearest prior demonstration that surfacing an
intermediate structured representation and closing an execution loop
materially improves text-to-CAD. It differs from our method in the two
respects that define our contribution: the expanded requirement carries no
per-field provenance (stated vs.\ assumed is not recorded), and its repair
signal is the program's own execution error and self-rendered views---not a
check against the user-confirmed specification. CADDesigner is also a single
general-purpose agent; we argue in \S\ref{sec:discussion:agents} that a small
decomposition with a deterministic kernel is easier to verify than one
monolithic agent, a point we return to when discussing the number of agents.

A second family validates programs against geometry through the CAD kernel
itself, but in the \emph{reverse-engineering} setting where a target mesh is
available: CADFit casts mesh-to-CAD recovery as volumetric-IoU-guided
optimization over executable programs, pruning operations whose marginal IoU
contribution is non-positive and nearly eliminating invalid programs
\cite{cadfit}; StepCAD conditions an LLM policy on the current intermediate
geometry and performs test-time tree search with local program edits under an
IoU objective \cite{stepcad}. Kernel-executed validation and per-step
geometry conditioning are close in spirit to our dynamic B-rep checks, but
their reference signal is the \emph{target mesh being reverse-engineered}---
ground truth that exists offline in reconstruction and never exists when
generating from a user's description.

Our verifier is anchored instead to the explicit set of user-confirmed
parameters and is performed in two deterministic channels: static alignment
of program literals to specification fields, and direct measurement of the
executed B-rep (radii, depths, feature counts, topology), with a vision
channel reserved for semantics the measurements cannot cover. Failures are
routed by cause: an unconfirmed assumption triggers another clarification
question, while a confirmed value the coder dropped triggers a
diagnostic-grounded regeneration rather than another question. This is closer
in spirit to specification- and property-based program verification
\cite{propertytesting} than to visual self-critique, and it requires no
ground-truth mesh at inference---the reference geometry that offline Chamfer
evaluation depends on \cite{chamfer_evals}.

\subsection{Positioning}\label{sec:related:positioning}

Table~\ref{tab:positioning} summarizes the axes that matter. No prior
text-to-CAD system combines intent clarification with a provenance-tracked
hand-off representation; no prior CAD program verifier is anchored to
user-confirmed requirements rather than to a target mesh or an unanchored
visual self-check; and no prior pipeline routes verification failures back to
clarification or regeneration by failure type.

\begin{table*}[t]
  \centering
  \caption{Positioning against the closest systems. ``GT at inference''
  denotes whether the method consumes a reference geometry at deployment
  time (not during offline evaluation).}
  \label{tab:positioning}
  \small
  \resizebox{\textwidth}{!}{%
  \begin{tabular}{lccccc}
    \toprule
    System & Clarification & Hand-off repr. & Deployment-time verification & Closed loop & GT at inference \\
    \midrule
    Text2CAD / DeepCAD \cite{text2cad,deepcad} & -- & command sequence & -- & -- & -- \\
    ProCAD \cite{procad} & one round, oracle & free text & -- & -- & no \\
    CADFusion / FutureCAD \cite{cadfusion,futurecad} & -- & programs + grounding & training-time & training loop & no \\
    CADCodeVerify / SEEK-CAD \cite{cadcodeverify,seekcad} & -- & code + renders & visual self-Q\&A & verify$\to$rewrite & no \\
    CADFit / StepCAD \cite{cadfit,stepcad} & -- & executable programs & kernel/IoU vs.\ target & search & \emph{yes (recon.\ input)} \\
    CAD-Editor / FlexCAD \cite{cadeditor,flexcad} & -- & editable text & -- & partial edits & -- \\
    Treellis \cite{treellis} & -- & assembly tree (humans) & -- & human edits & no \\
    \textbf{Ours} & \textbf{on-demand, multi-round} & \textbf{provenance spec} & \textbf{AST + B-rep vs.\ spec} & \textbf{by failure type} & \textbf{no} \\
    \bottomrule
  \end{tabular}%
  }
\end{table*}
